% Paper's fig:plaintext-mbqc, with labels for the classical and quantum outputs.
\begin{frame}{Plaintext MBQC}
  \centering
  \begin{tikzpicture}
    \path[use as bounding box] (-6.6,-3.15) rectangle (6.6,3.05);
    % Qcircuit's stick labels extend outside its measured matrix width.
    \node[inner sep=0pt,text=CVQCInk] at (-1.25,0.15) {%
      \scalebox{1.8}{%
        $\Qcircuit @C=1em @R=1em {
          \lstick{\ket{\psi}} & \ctrl{1} & \gate{HT^c} & \meter & \rstick{m}\cw \\
          \lstick{\ket{+}} & \control \qw & \qw & \qw & \rstick{X^mHT^c\ket{\psi}}\qw
        }$%
      }%
    };
  \end{tikzpicture}

  \note{
    \footnotesize
    This is the plaintext MBQC circuit from the manuscript's Technical
    Overview, Figure ``Gate teleportation.'' The input is psi on the upper
    wire and a fresh plus state on the lower wire. The two connected dots
    implement CZ. Apply $HT^c$ to the upper wire, then measure it in the
    computational basis. Equivalently, choose a measurement basis determined
    by c. Here $T=\operatorname{diag}(1,e^{i\pi/4})$ and c ranges from 0 to 7.
    If the outcome is m, the surviving lower wire holds
    $X^m HT^c\ket{\psi}$. This is the intended gate up to a known X
    correction; the upper output wire carries the classical measurement
    outcome. Repeating this gadget lets us implement gates through
    measurements, tracking corrections and adjusting later bases adaptively.
    Together with CZ, these operations are universal. The labels identify
    the classical outcome and the conditional state on the surviving wire.
    Explain the corrections in words; key updates and the encrypted version
    come later, when we introduce hiding and testing.
  }
\end{frame}
